\begin{textsample}{NEG Dim 4 – human – Score -3.00 – rj/rj\_inf\_e\_ef\_intro\_5\_p8.txt}  \label{ex:f4_neg_001}
Há de se considerar alguns princípios para a construção do ser humano e sociedade que queremos.
As \textbf{diretrizes} \textbf{curriculares} \textbf{nacionais} em relação às propostas pedagógicas devem garantir o cumprimento dos direitos do sujeito que perpassam todas as etapas de ensino e que sejam considerados princípios éticos, políticos, estéticos, cujas práticas pedagógicas norteadas por estes princípios, oportunizem à escola, a formação de sujeitos éticos, íntegros e comprometidos com o bem comum. É preciso garantir que todos os estudantes tenham direito a uma educação de \textbf{qualidade}, sem preconceitos de origem, raça, sexo, cor, idade e quaisquer outras formas de discriminação, conforme disposto no art. 3o, IV da Constituição Federal conjugado com o sentimento de pertencimento e da garantia da permanência nas escolas em todas as etapas, onde se valorize o sujeito de deveres e direitos comuns a todos, inclusive na reflexão sobre o direito de aprender.

% matched lemmas: curricular, diretriz, nacional, qualidade
\end{textsample}
